\chapter{EXPECTED OUTCOMES}

The primary deliverable of MotionBridge is a functioning real-time software system that accepts input from three synchronized USB webcams and produces two simultaneous outputs: a live PyBullet simulation window showing the simulated arm replicating the human subject's arm movement, and a continuously written CSV dataset file containing per-frame 3D joint coordinates and joint angles.

The expected dataset will contain full arm joint data (shoulder, elbow, wrist) and 21-point finger joint data per hand at the capture frame rate (target: 25--30~fps), with 3D joint position accuracy within acceptable bounds for downstream imitation learning pipelines. The system is expected to maintain real-time performance on standard laptop hardware without GPU acceleration.

Key performance metrics to be evaluated include: mean 3D joint reconstruction error (compared against a ground-truth reference pose), frames-per-second throughput of the full pipeline, simulation replication latency (time between human movement and simulated arm response), and dataset completeness (percentage of frames with all target joints detected across all cameras).

The dataset generated by the system will be structured and documented for public or academic release, enabling other researchers and practitioners to use it directly for robotic arm training without requiring their own motion capture infrastructure. This positions MotionBridge as both a working system and a data contribution to the robotics research community.
